\begingroup
\renewcommand{\qD}{q_{\mathbb D}^{0}}
\section{Foundation F51: exact disk Hessian and Bessel--Schur normalization}
\label{module:F51}
\noindent\textit{Audited source: \nolinkurl{convex_largeN_polya_szego_stage51_exact_disk_normalization.tex}.}
\subsection{Fourier and energy conventions}

For a real $2\pi$-periodic function $f$, use
\begin{equation}
 \widehat f_n=\frac1{2\pi}\int_0^{2\pi}f(\theta)e^{-in\theta}\,d\theta,
 \qquad
 f(\theta)=\sum_{n\in\mathbb Z}\widehat f_ne^{in\theta}.
\end{equation}
Our normalized circle inner product is
\begin{equation}\label{F51:eq:L2-star}
 \langle f,g\rangle_{L^2_*}
 :=\frac1{2\pi}\int_0^{2\pi}f(\theta)\overline{g(\theta)}\,d\theta
 =\sum_{n\in\mathbb Z}\widehat f_n\overline{\widehat g_n}.
\end{equation}
Let
\begin{equation}
 H_{\rm ng}^{s}
 :=\left\{f\in H^s(\Sone):
 \widehat f_0=\widehat f_1=\widehat f_{-1}=0\right\},
 \qquad
 \PiNG f=\sum_{|n|\ge2}\widehat f_ne^{in\theta}.
\end{equation}
The normalized homogeneous and inhomogeneous half-order forms are
\begin{align}
 \|f\|_{\dot H^{1/2}_*}^2
 &:=\sum_{|n|\ge2}|n|\,|\widehat f_n|^2,\label{F51:eq:hhalf}\\
 \qplus(f,g)
 &:=\sum_{|n|\ge2}(|n|+1)\widehat f_n\overline{\widehat g_n}.
 \label{F51:eq:qplus}
\end{align}
If $U$ is the decaying harmonic extension of $f$ to the half-cylinder
$\Sone\times(0,\infty)$, Parseval's identity gives
\begin{equation}\label{F51:eq:cylinder}
 \int_{\Sone\times(0,\infty)}|\nabla U|^2\,d\theta\,dy
 =2\pi\|f\|_{\dot H^{1/2}_*}^2.
\end{equation}
Thus a strip computation made in cylinder energy must be divided by $2\pi$
before it is inserted into \eqref{F51:eq:hhalf}.  This is the source of the two
constants
\begin{equation}\label{F51:eq:twoK}
 K_{\rm cyl}=\frac{\zeta(3)}{\pi^3},
 \qquad
 K_*:=\frac{K_{\rm cyl}}{2\pi}
 =\frac{\zeta(3)}{2\pi^4}.
\end{equation}
They are not competing normalizations of the same displayed form.

\subsection{The exact fan/support affine class}

Let
\begin{equation}
 \alpha_i>0,\qquad \sum_{i=1}^N\alpha_i=2\pi,
 \qquad \theta_{i+1}-\theta_i=\alpha_i
\end{equation}
cyclically, and put $\nu_i=(\cos\theta_i,\sin\theta_i)$.  On the edge cell
\begin{equation}
 K_i=[-\alpha_{i-1}/2,\alpha_i/2]
\end{equation}
write
\begin{equation}
 a=\alpha_{i-1},\quad b=\alpha_i,\quad
 s=\frac{x+a/2}{(a+b)/2},\quad
 r_- =\frac{z_i-z_{i-1}}a,\quad
 r_+=\frac{z_{i+1}-z_i}b.
\end{equation}

\begin{definition}[Exact trigonometric support trace]\label{F51:def:T}
For $z\in\R^N$, define $T_\alpha z$ on $K_i$ by
\begin{equation}\label{F51:eq:T}
 (T_\alpha z)_i(x)
 =z_i\cos x+\bigl((1-s)a_i^-+sa_i^+\bigr)\sin x,
\end{equation}
where
\begin{equation}\label{F51:eq:apm}
 a_i^-=r_-\frac a{\sin a}-z_i\tan\frac a2,
 \qquad
 a_i^+=r_+\frac b{\sin b}+z_i\tan\frac b2.
\end{equation}
The formulas on adjacent cells agree at their common endpoint, so
$T_\alpha z$ is continuous and belongs to $H^1(\Sone)$.
\end{definition}

Indeed, at the right endpoint of $K_i$, direct substitution gives
\begin{equation}
 (T_\alpha z)_i(b/2)
 =\frac{z_i+z_{i+1}}{2\cos(b/2)}
 =(T_\alpha z)_{i+1}(-b/2).
\end{equation}
The exact area row is
\begin{equation}\label{F51:eq:area-row}
 L_\alpha z=\sum_{i=1}^N\ell_i z_i,
 \qquad
 \ell_i=\tan\frac{\alpha_{i-1}}2+\tan\frac{\alpha_i}2,
\end{equation}
and the admissible coefficient space is $Z_\alpha=\ker L_\alpha$.  Let
\begin{equation}\label{F51:eq:V}
 V_\alpha:=\PiNG T_\alpha Z_\alpha\subset H_{\rm ng}^{1/2}.
\end{equation}

Let $B_\alpha^{\rm ex}\in H^{1/2}(\Sone)$ be the exact canonical fan
profile in this fixed material convention.  The normalization of that
profile is part of the definition; no factor is silently absorbed into it.
Set
\begin{equation}\label{F51:eq:b}
 b_\alpha:=\PiNG B_\alpha^{\rm ex},
 \qquad
 \mathcal A_\alpha=b_\alpha+V_\alpha.
\end{equation}
All endpoint values below depend only on the affine subspace
$\mathcal A_\alpha$, so replacing $B_\alpha^{\rm ex}$ by
$B_\alpha^{\rm ex}+T_\alpha z_0$ with $z_0\in Z_\alpha$ changes no value.

\subsubsection{Translations}

\begin{lemma}[Exact translation identity]\label{F51:lem:translations}
For $c\in\R^2$, put $z_i=c\cdot\nu_i$.  Then $L_\alpha z=0$ and
\begin{equation}\label{F51:eq:translation}
 T_\alpha z(\theta)=c\cdot(\cos\theta,\sin\theta).
\end{equation}
Consequently $\PiNG T_\alpha z=0$.
\end{lemma}

\begin{proof}
Let $\tau_i=(-\sin\theta_i,\cos\theta_i)$.  From
\begin{equation}
 \nu_{i-1}=\nu_i\cos a-\tau_i\sin a,
 \qquad
 \nu_{i+1}=\nu_i\cos b+\tau_i\sin b
\end{equation}
one obtains from \eqref{F51:eq:apm}
\begin{equation}
 a_i^-=a_i^+=c\cdot\tau_i.
\end{equation}
Equation \eqref{F51:eq:T} is therefore exactly
\begin{equation}
 (c\cdot\nu_i)\cos x+(c\cdot\tau_i)\sin x
 =c\cdot\nu(\theta_i+x),
\end{equation}
which proves \eqref{F51:eq:translation}.  The area row vanishes because
translations preserve area.  Algebraically, it is the closure identity
$\sum_i\ell_i\nu_i=0$ for the tangential polygon with normal fan $\alpha$.
The last assertion follows from the Fourier support of
$c\cdot\nu(\theta)$.
\end{proof}

Thus centering only chooses a representative of the translation orbit.  It
does not add a finite-rank loss to the nongeometric endpoint Schur value.

\subsection{The disk multiplier}

Let $j=j_{0,1}$ be the first positive zero of $J_0$.  For $n\ge2$, set
\begin{equation}\label{F51:eq:kn}
 k_n=j\frac{J_{n+1}(j)}{J_n(j)}.
\end{equation}

\begin{lemma}[A rational enclosure for $j$]\label{F51:lem:j-bound}
One has
\begin{equation}\label{F51:eq:j-bound}
 2.4<j<2.405,
 \qquad \frac{j^2}{4}<1.447.
\end{equation}
\end{lemma}

\begin{proof}
The series
\begin{equation}
 J_0(x)=\sum_{m\ge0}\frac{(-1)^m(x^2/4)^m}{(m!)^2}
\end{equation}
is alternating with decreasing terms on the indicated interval after its
first term.  Rational partial sums through orders eight and nine give
\begin{align*}
 0.0025076832943&<J_0(12/5)<0.0025076834966,\\
 -0.000090558004&<J_0(481/200)<-0.000090557793.
\end{align*}
Moreover $J_0'=-J_1$, and
\begin{equation}
 J_1(x)=\frac x2\sum_{m\ge0}
 \frac{(-1)^m(x^2/4)^m}{m!(m+1)!}>0
 \qquad(0<x\le2.405),
\end{equation}
because the alternating sum is at least $1-(2.405)^2/8>0$.
Hence $J_0$ is strictly decreasing there and its first zero lies in the
displayed interval.  The final bound follows by squaring $2.405$.
\end{proof}

\begin{theorem}[Positive order-$(-1)$ splitting]\label{F51:thm:bessel}
For every $n\ge2$,
\begin{equation}\label{F51:eq:recurrence}
 1+j\frac{J_n'(j)}{J_n(j)}=n+1-k_n,
\end{equation}
and
\begin{equation}\label{F51:eq:k-bounds}
 0<k_n\le \frac{2x}{n+1-x}<2,
 \qquad k_n<\frac4n,
 \qquad x:=\frac{j^2}{4}<1.447.
\end{equation}
In particular
\begin{equation}\label{F51:eq:weight-equivalence}
 \frac13(n+1)\le n+1-k_n\le n+1
 \qquad(n\ge2).
\end{equation}
\end{theorem}

\begin{proof}
The Bessel recurrence
\begin{equation}
 jJ_n'(j)=nJ_n(j)-jJ_{n+1}(j)
\end{equation}
gives \eqref{F51:eq:recurrence}.  Write
\begin{equation}
 J_n(j)=\frac{(j/2)^n}{n!}F_n,
 \qquad
 F_n=\sum_{m\ge0}\frac{(-1)^mx^m}{m!(n+1)_m}.
\end{equation}
For $n\ge2$, the ratio of consecutive absolute terms is at most
$x/(n+1)<1/2$.  Alternating-series bounds give
\begin{equation}
 0<1-\frac{x}{n+1}\le F_n\le1.
\end{equation}
Consequently
\begin{equation}
 0<k_n=\frac{2x}{n+1}\frac{F_{n+1}}{F_n}
 \le\frac{2x}{n+1-x}.
\end{equation}
The right-hand side is decreasing in $n$ and is less than $1.864$ at
$n=2$.  Direct cross multiplication, using $x<1.447$, also gives
$2x/(n+1-x)<4/n$ for every $n\ge2$.  Hence
$n+1-k_n\ge n-1\ge(n+1)/3$, proving
\eqref{F51:eq:weight-equivalence}.
\end{proof}

Define the positive self-adjoint multiplier
\begin{equation}
 K_{-1}e^{in\theta}=k_{|n|}e^{in\theta},
 \qquad |n|\ge2,
\end{equation}
and the reduced disk form
\begin{equation}\label{F51:eq:qD}
 \qD(f,g):=\qplus(f,g)-k(f,g),
 \qquad
 k(f,g):=\langle f,K_{-1}g\rangle_{L^2_*}.
\end{equation}
The physical disk quadratic Taylor form (equivalently, one half of the
second variation of the scale-invariant deficit at the disk) is
\begin{equation}\label{F51:eq:physical-Q}
 \mathsf Q_D(f,g):=2j^2\qD(f,g).
\end{equation}
Thus the full Hessian is $2\mathsf Q_D=4j^2\qD$.  The factor $2j^2$ is
not included in either $\qplus$ or $\qD$.

\begin{corollary}[Coercivity and negative order]\label{F51:cor:coercivity}
For $f\in H_{\rm ng}^{1/2}$,
\begin{equation}\label{F51:eq:coercivity}
 \frac13\qplus(f,f)\le\qD(f,f)\le\qplus(f,f).
\end{equation}
Moreover
\begin{equation}\label{F51:eq:minus}
 |k(f,g)|\le C_j\|f\|_{H^{-1/2}}\|g\|_{H^{-1/2}}.
\end{equation}
\end{corollary}

\begin{proof}
Equation \eqref{F51:eq:coercivity} is \eqref{F51:eq:weight-equivalence} summed over
Fourier modes.  By \eqref{F51:eq:k-bounds}, $k_n\le C_j/n$; Cauchy--Schwarz with
weight $1/n$ gives \eqref{F51:eq:minus}.
\end{proof}

\subsection{Exact terminal Schur identity}

For the affine class \eqref{F51:eq:b}, define
\begin{align}
 \Pplus(\alpha)
 &:=\inf_{v\in V_\alpha}\qplus(b_\alpha+v,b_\alpha+v),\label{F51:eq:Pplus}\\
 \widetilde{\PD}(\alpha)
 &:=\inf_{v\in V_\alpha}\qD(b_\alpha+v,b_\alpha+v),\label{F51:eq:PDtilde}\\
 \PD(\alpha)&:=2j^2\widetilde{\PD}(\alpha).
 \label{F51:eq:PD}
\end{align}
Let $f_\alpha\in\mathcal A_\alpha$ be the unique $\qplus$-minimal
residual.  Thus
\begin{equation}\label{F51:eq:galerkin}
 \qplus(f_\alpha,v)=0\qquad(v\in V_\alpha).
\end{equation}
On $V_\alpha$, define
\begin{equation}
 A_{D,\alpha}[v,w]=\qD(v,w),
 \qquad g_\alpha(v)=k(f_\alpha,v).
\end{equation}

\begin{theorem}[Terminal Bessel--Schur formula]\label{F51:thm:schur}
For every positive fan,
\begin{equation}\label{F51:eq:schur}
 \widetilde{\PD}(\alpha)
 =\Pplus(\alpha)-k(f_\alpha,f_\alpha)
 -\langle g_\alpha,A_{D,\alpha}^{-1}g_\alpha\rangle.
\end{equation}
The inverse and dual pairing are taken on the quotient of $V_\alpha$ by its
zero trace; by \eqref{F51:eq:coercivity} they are well defined.
\end{theorem}

\begin{proof}
Every element of $\mathcal A_\alpha$ has a unique representation
$f_\alpha+w$ modulo the zero trace, with $w\in V_\alpha$.  Using
\eqref{F51:eq:galerkin},
\begin{align*}
 \qD(f_\alpha+w,f_\alpha+w)
 ={}&\qplus(f_\alpha,f_\alpha)-k(f_\alpha,f_\alpha)\\
 &+A_{D,\alpha}[w,w]-2g_\alpha(w).
\end{align*}
Completing the square in the coercive form $A_{D,\alpha}$ gives
\eqref{F51:eq:schur}.
\end{proof}

\subsection{The uniform fan}

Let $\alpha_*=a\mathbf1$, $a=2\pi/N$.  Rotation by $a$ acts on coefficient
vectors by the cyclic shift and on traces by the unitary rotation
$R_af(\theta)=f(\theta-a)$.

\begin{lemma}[Uniform Bessel forcing vanishes]\label{F51:lem:gzero}
Assume the canonical choice of $B_{\alpha_*}^{\rm ex}$ is cyclically
equivariant.  Then
\begin{equation}\label{F51:eq:gzero}
 g_{\alpha_*}=0.
\end{equation}
\end{lemma}

\begin{proof}
The affine class, $\qplus$, and the canonical datum are invariant under
$R_a$, so uniqueness of the $\qplus$ projection gives
$R_af_{\alpha_*}=f_{\alpha_*}$.  The multiplier $K_{-1}$ commutes with
$R_a$.  Hence for $z\in Z_{\alpha_*}$,
\begin{equation}
 k(f_{\alpha_*},\PiNG T_{\alpha_*}z)
 =k\left(f_{\alpha_*},\PiNG T_{\alpha_*}
 \frac1N\sum_{m=0}^{N-1}\sigma^m z\right),
\end{equation}
where $\sigma$ is the cyclic shift.  At the uniform fan,
$L_{\alpha_*}z=0$ is equivalent to $\sum_i z_i=0$, so the averaged
coefficient vector is zero.  Therefore the displayed pairing vanishes for
every element of $V_{\alpha_*}$.
\end{proof}

Combining Theorem~\ref{F51:thm:schur} and Lemma~\ref{F51:lem:gzero} gives the exact
endpoint difference
\begin{align}\label{F51:eq:endpoint-difference}
 \widetilde{\PD}(\alpha)-\widetilde{\PD}(\alpha_*)
 ={}&\Pplus(\alpha)-\Pplus(\alpha_*)\\
 &-\bigl[k(f_\alpha,f_\alpha)
          -k(f_{\alpha_*},f_{\alpha_*})\bigr]\notag\\
 &-\langle g_\alpha,A_{D,\alpha}^{-1}g_\alpha\rangle.\notag
\end{align}
\endgroup
